\documentclass[11pt]{article}
\usepackage{amsmath,amssymb,amsthm}
\usepackage{graphicx}
\usepackage[margin=1in]{geometry}

\newtheorem{definition}{Definition}
\newtheorem{theorem}{Theorem}
\newtheorem{example}{Example}
\newtheorem{remark}{Remark}

\title{Lagrangian Mechanics and the Principle of Least Action}
\author{PHYS 6801 --- Classical Mechanics (Example)}
\date{January 15, 2026}

\begin{document}
\maketitle

\section{Principle of Least Action}

% block_001 | content_type: definition | confidence: high
\begin{definition}[Action]
The \emph{action} $S$ associated with a path $q(t)$ is defined as the time integral of the
Lagrangian $L$ along that path:
\begin{equation}
  S = \int L\bigl(q(t), \dot{q}(t), t\bigr)\, dt.
  \label{eq:action}
\end{equation}
\end{definition}

% block_002 | content_type: derivation | confidence: high | depends_on: block_001
We seek paths for which the action is stationary under variations $\delta q(t)$ of the trajectory.
Requiring the first variation of the action to vanish gives the condition
\begin{equation}
  \delta S = 0.
  \label{eq:stationary-action}
\end{equation}
This variational principle selects the physical path among all paths with fixed endpoints.

\section{Euler--Lagrange Equation}

% block_003 | content_type: equation | confidence: medium
% REVIEW: Lower portion of equation slightly smudged; verify partial derivative notation.
Imposing \eqref{eq:stationary-action} on \eqref{eq:action} yields the Euler--Lagrange equation:
\begin{equation}
  \frac{d}{dt}\left(\frac{\partial L}{\partial \dot{q}}\right)
  - \frac{\partial L}{\partial q} = 0.
  \label{eq:euler-lagrange}
\end{equation}

\end{document}
